\documentclass[conference]{IEEEtran}

% -------------------------------------------------------
%  Packages
% -------------------------------------------------------
\usepackage{cite}
\usepackage{amsmath,amssymb,amsfonts}
\usepackage{graphicx}
\usepackage{textcomp}
\usepackage{xcolor}
\usepackage{booktabs}
\usepackage{hyperref}
\usepackage{algorithmic}
\usepackage{array}
\usepackage{balance}
\usepackage{float}
\usepackage{multirow}

\hypersetup{
    colorlinks=true,
    linkcolor=blue,
    citecolor=blue,
    urlcolor=blue
}

% -------------------------------------------------------
%  Title
% -------------------------------------------------------
\title{Improve performance of existing online trackers for small object detection [UAV videos] \\
\textbf{\large Team: The Convolutioners}

}

% -------------------------------------------------------
%  Authors
% -------------------------------------------------------
\author{

\IEEEauthorblockN{Aagam Sheth}
\IEEEauthorblockA{
AU2544006\\
M.Tech (CSE)
}
\and
\IEEEauthorblockN{Mahima Parekh}
\IEEEauthorblockA{
AU2544003\\
M.Tech (CSE)
}
\and
\IEEEauthorblockN{Khevan Thanki}
\IEEEauthorblockA{
AU23L20002\\
BS (Hons) CS
}
\and
\IEEEauthorblockN{Riddhi Bhargava}
\IEEEauthorblockA{
AU22220028\\
BS (Hons) CS
}

}

% -------------------------------------------------------
\begin{document}
\maketitle

% -------------------------------------------------------
%  Abstract
% -------------------------------------------------------
\begin{abstract}
Small object detection (SOD) and multi-object tracking (MOT) in unmanned aerial vehicle (UAV) videos have emerged as significant research areas in computer vision. Identification of small objects in various environments has proven challenging due to factors such as small object sizes, low frame resolutions, and dynamic scene conditions, leading to missed detections and increased false positives in current object detectors. Therefore, this project proposes using a deep learning framework for robust SOD and MOT in UAV videos, incorporating SAHI, a tiling strategy framework, to enhance the performance of the You Only Look Once (YOLO) detector on the VisDrone2019-MOT dataset. A custom tracking module is developed that improves re-identification by integrating a Siamese Network Module and a Kalman filter for motion prediction. This project aims to provide a modular framework for evaluating tracking performance across diverse UAV scenarios and to evaluate the proposed system against established benchmarks.
\end{abstract}

% -------------------------------------------------------
%  Keywords
% -------------------------------------------------------
\begin{IEEEkeywords}
Small Object Detection, Multi-Object Detection, Multi-Scale Learning, UAV, VisDrone, Kalman filter.
\end{IEEEkeywords}

% -------------------------------------------------------
%  I. Introduction
% -------------------------------------------------------
\section{Introduction}
Accurate detection and classification of objects in unmanned aerial vehicles (UAVs) videos are critical for numerous real-world applications, including traffic monitoring, urban planning, and surveillance \cite{ref1}. However, current object detection methods often exhibit suboptimal performance when applied to UAV scenarios, primarily due to the distinct challenges introduced by aerial viewpoints and imaging conditions.\\
Unlike objects in proper scales, small object detection (SOD) is much more challenging due to the extremely small size and low signal-to-noise ratio in aerial images \cite{ref2}. Detectors such as YOLO, while highly efficient for real-time object detection, often struggle with SOD in UAV scenarios due to their limited receptive field, leading to missed detections and false positives \cite{ref3}.
Although existing SOD algorithms have enhanced detection performance, the detection accuracy in practical applications still fails to meet requirements and leaves substantial room for improvement \cite{ref4}.
Therefore, this project proposes a unified deep learning framework for robust small object detection and multi-object tracking (MOT) in UAV videos. The proposed system integrates a tiling-based preprocessing strategy with a YOLO11 detector trained on the VisDrone2019-MOT dataset. \\
To address the limitations of existing trackers in aerial scenarios, a custom tracking module is developed based on the BoT-SORT tracker while replacing its standard re-identification module with a Siamese network trained specifically on UAV object pairs. Motion continuity across frames is maintained through a Kalman filter, which predicts the next bounding box position and compensates for frames in which the detector fails to localize an object. The system is evaluated on the VisDrone2019-MOT benchmark using TrackEval metrics including MOTA, MOTP, and IDF1, with additional analysis of localization accuracy and false positive rates in cluttered multi-object scenes.

The remainder of this paper is structured as follows. Section~\ref{sec:method}
describes the proposed methodology. Section~\ref{sec:results} presents
preliminary results. Section~\ref{sec:discussion} discusses findings and
implications. Section~\ref{sec:conclusion} concludes the paper.

% -------------------------------------------------------
%  II. Methodology
% -------------------------------------------------------
\section{Methodology}
\label{sec:method}

The proposed system consists of five phases as illustrated in the overall
pipeline: data preparation, baseline evaluation, tiling strategy implementation on YOLO, designing of Siamese Network, and developing a custom tracker.
\subsection{Phase 1: Data Preparation}

The VisDrone2019-MOT dataset is used to train and test the designed architecture. The dataset contains ten object categories including pedestrian, bicycle, car, van, truck, and motor. \\
In the preparation stage, VisDrone annotations are transformed into YOLO format, which is required to train the YOLO11 detector, where each object's class ID and bounding box are normalized to the image dimensions as center coordinates (cx, cy, w, h). Annotations marked as ignored regions and boxes with invalid dimensions are filtered out to ensure label quality. The dataset is systematically divided into training and validation sets, on which YOLO11 is trained and evaluated respectively.
\\
Additionally, for the training data for the Siamese network in Phase 4, object crops are extracted from the VisDrone2019-MOT training sequences by localising each annotated bounding box and saving the corresponding image patch.

\subsection{Phase 2: Baseline Evaluation}
YOLOv8n was evaluated on the VisDrone dataset to establish baseline performance metrics. 
This baseline reflects the performance of a standard YOLOv8n detector on full-resolution UAV frames without any scale-aware preprocessing and serves as a reference point for subsequent improvements. 
Additionally, YOLO11 was fine-tuned on the VisDrone-MOT dataset to assess performance gains compared to the non-fine-tuned model.

\subsection{Phase 3: Fine-tuned YOLO11 + SAHI tiling strategy}
A key limitation of the baseline YOLO model is that full VisDrone frames are downsampled during training and inference, causing small objects to become effectively undetectable after resizing. To address this, a tiling-based preprocessing strategy was applied using the SAHI (Slicing Aided Hyper Inference) framework.

In this approach, each input frame is divided into overlapping crops of 640×640 pixels, and the fine-tuned YOLO11 detector is applied independently to each tile. Multiple detections may be produced for the same object across adjacent tiles due to the overlapping regions. This is resolved by applying Non-Maximum Suppression (NMS) to merge redundant bounding boxes and retain the most confident detections, thereby ensuring consistent and accurate object localization.
The tiling-augmented detector was evaluated on the same VisDrone2019-MOT validation split as the baseline, allowing a direct comparison of the evaluation metrics between the two approaches.

\subsection{Phase 4: Siamese Network for Re-Identification}
The pipeline begins with extracting ground-truth object crops from VisDrone annotations using a preprocessing script. Pairwise training samples are then constructed, generating positive pairs (same identity across frames) and negative pairs (different identities, including visually similar categories).

To address the limitations of standard ReID models, a Siamese network is employed. The network consists of two weight-sharing convolutional branches with a ResNet18 backbone, which learn an embedding function \( f(x) \) that maps each input crop to a normalized feature space.

Similarity between object pairs is computed using the Euclidean distance between embeddings, where smaller distances indicate higher likelihood of identity match. The network is trained using contrastive loss to minimize intra-class distance while enforcing a margin between inter-class pairs. During evaluation, a distance threshold is applied to classify pairs, and performance is assessed using standard classification metrics and distance distribution analysis, (Table IV).

\subsection{Phase 5: BoT-SORT Tracker Design}
The tracking pipeline combines motion prediction,spatial matching, and appearance based reidentification into a two-stage association strategy. At each frame, existing tracks are predicted forward using a constant-velocity Kalman filter, and incoming detections are split into high-confidence and low-confidence groups. In the first stage, high-confidence detections are matched to all active tracks using a blended cost of IoU distance and cosine appearance distance, resolved via the Hungarian algorithm. In the second stage, remaining unmatched tracks are associated with low-confidence detections using IoU cost alone. Stale tracks are pruned while unmatched detections that satisfy birth thresholds are initialised as new tracks.
The Siamese network is central to the appearance matching component. Each detection patch is cropped from the RGB frame, resized to 128×128, and passed through a single forward branch to produce a 128-dimensional L2-normalised embedding. Each active track maintains an appearance embedding updated via exponential moving average, which is used to compute cosine similarity during first-stage association. This is especially critical in UAV footage where small, closely spaced objects produce ambiguous IoU overlaps that geometry alone cannot resolve. YOLOv8 serves as the detection front-end, with the detector and tracker kept fully decoupled so that improvements in detection recall propagate directly into tracking performance.

% -------------------------------------------------------
%  III. Results
% -------------------------------------------------------

\section{Results}
\label{sec:results}

\subsection{Baseline Evaluation}

\begin{table}[H]
\centering
\caption{YOLO11 + SAHI performance with existing BoT-SORT}
\begin{tabular}{l c}
\toprule
\textbf{Metric} & \textbf{Value} \\
\midrule
MOTA & 17.28 \% \\
MOTP & 26.01 \% \\
IDF1 & 48.79 \% \\
Precision & 61.12 \% \\
Recall & 49.07 \% \\
Mostly Tracked (MT) & 224 \\
Mostly Lost (ML) & 271 \\
False Positives (FP) & 35629 \\
False Negatives (FN) & 58124 \\
ID Switches & 657 \\
FPS & 0.5 \\
\bottomrule
\end{tabular}
\label{tab:botsort_summary}
\end{table}

\subsection{Siamese Network for Re-Identification}

To enhance appearance-based data association in UAV tracking, we train a Siamese re-identification network on VisDrone object crops. The model learns an embedding function using contrastive loss, enforcing intra-class compactness while maintaining interclass separation. The learned embeddings are evaluated using classification metrics and distance-based analysis.

\begin{table}[H]
    \centering
    \caption{Epoch-wise Training Performance of the Siamese Network}
    \begin{tabular}{c c c c c c}
        \hline
        \textbf{Epoch} & \textbf{Loss} & \textbf{Accuracy} & \textbf{Precision} & \textbf{Recall} & \textbf{F1-score} \\
        \hline
        1 & 2.6657 & 0.6636 & 0.6032 & 0.9718 & 0.7444 \\
        2 & 0.2155 & 0.6771 & 0.6069 & 0.9986 & 0.7550 \\
        3 & 0.1390 & 0.7165 & 0.6347 & 1.0000 & 0.7765 \\
        4 & 0.0982 & 0.7914 & 0.7032 & 1.0000 & 0.8257 \\
        5 & 0.0829 & 0.7713 & 0.6866 & 1.0000 & 0.8142 \\
        \hline
    \end{tabular}
    \label{tab:epoch_results}
\end{table}



\begin{table}[H]
    \centering
    \caption{Embedding Distance Statistics Across Training}
    \begin{tabular}{l c c}
        \hline
        \textbf{Metric} & \textbf{Epoch 1} & \textbf{Epoch 5} \\
        \hline
        Same-Object Distance & 0.4849 & 0.2664 \\
        Different-Object Distance & 1.1573 & 1.0428 \\
        \hline
    \end{tabular}
    \label{tab:embedding_distances_horizontal}
\end{table}

\begin{figure*}[!t]
    \centering
    \includegraphics[width=\textwidth]{training_metrics.png}
    \caption{Training dynamics of the Siamese network. The plot illustrates rapid loss convergence and steady improvements in accuracy, precision, recall, and F1-score. The saturation of recall and stabilization of metrics after epoch 4 indicate convergence to a near-optimal solution.}
    \label{fig:training_metrics}
\end{figure*}


\subsection{BoT-SORT Tracker Results}
The system was evaluated on the VisDrone2019-MOT benchmark. The evaluated pipeline combines YOLOv8 detections with a Siamese-enhanced BoT-SORT tracker. The performance metrics obtained are summarized in Table~\ref{tab:visdrone_results}.

\begin{table}[H]
\centering
\caption{Performance on VisDrone2019-MOT Benchmark}
\begin{tabular}{c c}
\hline
\textbf{Metric} & \textbf{Value} \\
\hline
MOTA & 0.129 \\
IDF1 & 0.224 \\
MOTP & 0.222 \\
Mean HOTA@0.5 & 0.260 \\
Mean DetA@0.5 & 0.142 \\
Mean AssA@0.5 & 0.504 \\
IDP & 0.817 \\
IDR & 0.130 \\
Overall FPS & 12.69 \\
\hline
\end{tabular}
\label{tab:visdrone_results}
\end{table}






% -------------------------------------------------------
%  IV. Discussion
% -------------------------------------------------------
\section{Discussion}
\label{sec:discussion}

Table~\ref{tab:botsort_summary} shows that fine-tuning YOLO11 with SAHI tiling improves detection coverage, especially for small and dense objects. This leads to better recall and MOTA, though an increase in false positives is also observed due to overlapping tiles,  suggesting need for better post-processing.

Table~\ref{tab:epoch_results} shows the Siamese model converges fast, with best performance around epoch 4 before slight over fitting. Recall reaches 1.0 early,which helps capture same object pairs but also makes the decision boundary bit permissive. Table~\ref{tab:embedding_distances_horizontal} confirms good embedding structure, with reduced same object distances and clear separation from different object pairs.

Figure~\ref{fig:training_metrics} shows stable training, with sharp loss drop and metric saturation. The model focuses more on reducing false negatives first, and later improves precision.

For tracking, the Siamese + BoTSORT setup shows strong association (AssA = 0.504, IDP = 0.817) but still suffers from low recall (IDR = 0.130), leading to modest overall scores (MOTA = 0.129, HOTA = 0.260). Overall, results suggest that while association is strong, improving detection remains critical for better performance.

% -------------------------------------------------------
%  V. Conclusion
% -------------------------------------------------------
\section{Conclusion}
\label{sec:conclusion}

This study presents a modular framework for improving small object detection and multi-object tracking in UAV videos. By integrating a SAHI based tiling strategy with a fine-tuned YOLO detector, the system enhances small object detection. A Siamese network is employed  to improve re-identification, while a BoT-SORT-based tracker with Kalman filtering ensures temporal consistency. Although the full integration of the YOLO + tiling module with the BoTSORT Siamese tracker remains ongoing, preliminary results are promising. The findings indicate that combining scale-aware detection with appearance-based tracking can improve performance, with further integration and optimization expected to yield more robust and accurate tracking outcomes.



% -------------------------------------------------------
%  References
% -------------------------------------------------------
\balance
\begin{thebibliography}{00}
\bibitem{ref1}
Xiao, Y. et al. (2025) ‘UAV video vehicle detection: Benchmark and baseline’, IEEE Transactions on Geoscience and Remote Sensing, 63, pp. 1–14. doi:10.1109/tgrs.2025.3534524.

\bibitem{ref2}
Wang, J. et al. (2021) ‘Tiny object detection in aerial images’, 2020 25th International Conference on Pattern Recognition (ICPR), pp. 3791–3798. doi:10.1109/icpr48806.2021.9413340. 

\bibitem{ref3}
Liu, M. et al. (2020) ‘UAV-YOLO: Small object detection on unmanned aerial vehicle perspective’, Sensors, 20(8), p. 2238. doi:10.3390/s20082238. 

\bibitem{ref4}
Feng, F. et al. (2024) ‘Improved yolov8 algorithms for small object detection in aerial imagery’, Journal of King Saud University - Computer and Information Sciences, 36(6), p. 102113. doi:10.1016/j.jksuci.2024.102113.

\end{thebibliography}

\end{document}